\section{Protection and Bench Operation}
\label{sec:req-safety}

Controlled energization and rated measurement equipment were used during testing. The high-voltage assembly was checked for remaining charge before handling. Section~\ref{sec:method-safety} describes the bench arrangement. The prototype was intended for supervised laboratory testing.

The V1.1 drawing includes an input fuse, surge suppression, line filtering, and a CAP300DG discharge circuit, along with protection components in the auxiliary and USB-C sections. Chapter~\ref{ch:hardware} explains their roles. These functions support laboratory operation; qualification still requires tests of the complete assembly.

Precharge applies limited excitation to power the isolated feedback circuit. It ends when the feedback threshold or configured attempt limit is reached. High-impedance classification lowers the burst limit and voltage window. Eligible light/no-load overvoltage invokes soft holdoff within the hard margin; other overvoltage conditions latch overvoltage protection (OVP) until the defined clear or reset action. Burst enable (PA10) is held low to inhibit further excitation, without discharging stored energy. Chapter~\ref{ch:control} and Appendix~\ref{app:equations} specify the conditions.

Primary half-bridge-current feedback (PA2) supports load classification and diagnostics. The configurable current trip is disabled by default, and the active voltage loop does not call legacy foldback helpers. Any current-limited test must identify the active limiter: output controller, AC source, electronic load, or explicitly enabled firmware threshold.

Insulation qualification, leakage-current, abnormal-fault, surge/ESD, and thermal-protection tests remain unfinished. The completed checks cover supervised laboratory operation, not product qualification.
